% docs/manual/operations/appendices/B-remote-config.tex
% 附录 B：SSH / Remote 配置参考

\chapter{SSH / Remote 配置参考}

完整 SSH 配置参考。详细背景见第 3 章。

\section{Connection profile schema}

\file{\~{}/.openbench\_wizard/connections.yaml}：

\begin{minted}{yaml}
connections:
  <profile_name>:
    host: <user@hostname>      # 必填
    port: 22                    # 默认 22
    key_path: ~/.ssh/<keyname>  # 私钥路径；推荐
    password: null              # 不推荐；用密钥
    remote_workdir: <path>      # 远程工作目录
    auth_method: key | password # 认证方式
\end{minted}

\section{凭据存储路径}

\begin{longtable}{p{3.5cm} p{9cm}}
  \toprule
  系统 & 路径 \\
  \midrule
  \endhead
  macOS & \texttt{\~{}/.openbench\_wizard/} \\
  Linux & \texttt{\~{}/.openbench\_wizard/} \\
  Windows & \texttt{\%USERPROFILE\%\textbackslash{}.openbench\_wizard\textbackslash{}} \\
  \bottomrule
\end{longtable}

文件清单：

\begin{itemize}
  \item \texttt{connections.yaml}：profile 列表（明文）
  \item \texttt{credentials.json}：加密的密码（如有）
  \item \texttt{encryption\_salt.bin}：fernet 派生用 salt
\end{itemize}

\section{权限要求}

\begin{minted}{bash}
chmod 600 ~/.openbench_wizard/credentials.json
chmod 600 ~/.openbench_wizard/encryption_salt.bin
chmod 600 ~/.openbench_wizard/connections.yaml
chmod 700 ~/.openbench_wizard/
\end{minted}

如系统配 umask 不默认 077，每次写入后手动确认权限。

\section{ssh-agent 集成}

\openbench{} GUI 使用 paramiko，会读 \texttt{\~{}/.ssh/config} 与 \texttt{\$SSH\_AUTH\_SOCK}。如使用 ssh-agent：

\begin{minted}{bash}
eval "$(ssh-agent)"
ssh-add ~/.ssh/hpc_key
# 之后 GUI 不需输入 passphrase
\end{minted}

\section{堡垒机 / ProxyJump}

\openbench{} GUI 直接读 \texttt{\~{}/.ssh/config}，标准 ProxyJump 自动生效：

\begin{minted}{text}
# ~/.ssh/config
Host bastion
  HostName bastion.example.edu
  User alice

Host hpc01
  HostName hpc01.internal
  User alice
  ProxyJump bastion
  IdentityFile ~/.ssh/hpc_key
\end{minted}

GUI profile 中 \texttt{host} 设为 \texttt{hpc01}（用 SSH config 别名），不需要重复 ProxyJump 配置。

\section{端口转发场景}

如要通过 SSH 转发 HPC 上的 Jupyter / TensorBoard / Prometheus：

\begin{minted}{bash}
ssh -L 8888:localhost:8888 hpc01
# 之后本地 http://localhost:8888 访问远程服务
\end{minted}

\openbench{} 不直接管端口转发，但远程 \cli{run} 完成后用此方式查看远程报告：

\begin{minted}{bash}
# 远程
ssh hpc01 "cd ~/eval && python -m http.server 8000"

# 本地（另开窗口）
ssh -L 8000:localhost:8000 hpc01
# 浏览器访问 http://localhost:8000/output/<run>/reports/report.html
\end{minted}

\section{已知问题速查}

\begin{longtable}{p{4cm} p{8.5cm}}
  \toprule
  症状 & 处置 \\
  \midrule
  \endhead
  Permission denied (publickey) & \texttt{chmod 600 \~{}/.ssh/key}；确认公钥在远程 authorized\_keys \\
  Host key verification failed & \texttt{ssh-keygen -R hostname}，重新接受 \\
  Connection timed out & VPN / 防火墙；\texttt{nc -vz host port} 测连通 \\
  paramiko bad packet & 服务器 SSH 版本太老；升级或用 system ssh \\
  GUI Test 红但 ssh 命令行成功 & GUI 没读 \texttt{\~{}/.ssh/config}（罕见），把 host 写完整 \\
  \bottomrule
\end{longtable}
